\documentclass[11pt,a4paper]{article}
\usepackage[T1]{fontenc}
\usepackage{lmodern}
\usepackage[margin=26mm]{geometry}
\usepackage{amsmath,amssymb,amsthm,booktabs,array,longtable}
\usepackage{microtype,multicol}
\usepackage[hidelinks]{hyperref}
\setlength{\parskip}{3pt}
\setlength{\emergencystretch}{2em}
\raggedbottom
\allowdisplaybreaks[1]
\clubpenalty=10000
\widowpenalty=10000
\newtheorem{theorem}{Theorem}[section]
\newtheorem{lemma}[theorem]{Lemma}
\newtheorem{corollary}[theorem]{Corollary}
\theoremstyle{remark}
\newtheorem{remark}[theorem]{Remark}
\newcommand{\Z}{\mathbb Z}
\newcommand{\Q}{\mathbb Q}
\newcommand{\N}{\mathbb N}
\newcommand{\F}{\mathbb F}
\newcommand{\ps}[1]{#1[[x]]}
\newcommand{\Rev}{\operatorname{Rev}}
\newcommand{\rev}{^{\langle-1\rangle}}
\newcommand{\Plus}{\operatorname{Plus}}
\newcommand{\Minus}{\operatorname{Minus}}
\newcommand{\coeff}[2]{[x^{#1}]#2}
\newcommand{\E}{\mathcal E}
\Urlmuskip=0mu plus 1mu
\hypersetup{pdftitle={A393172: A Catalan substitution congruence modulo four},pdfauthor={Dev Goyal},pdfsubject={Separate manuscript with AI disclosure and Lean proof map}}
\title{A393172\\[4pt]\large A Catalan substitution congruence modulo four}
\author{Dev Goyal}
\date{24 October 2026}
\begin{document}
\maketitle
\begin{center}
\fbox{\begin{minipage}{0.9\linewidth}\small
\textbf{AI disclosure.} ChatGPT (OpenAI) generated the mathematical arguments,
Lean proofs, and this manuscript in response to Dev Goyal's requests.
Dev Goyal is the sole human author. The source sequence and conjecture or
observation are credited to Paul D. Hanna.
\textbf{Independent human review is pending.}
\end{minipage}}
\end{center}
\begin{abstract}
The unique normalized integral solution of $A(x-x^2)=x+A^2$ satisfies $A\equiv2C-x\pmod4$, where $C=x+C^2$. Hence its coefficients above degree one have residue two precisely at powers of two, and zero otherwise. The comparison satisfies the literal source equation.
\end{abstract}

\section{Source, notation, and scope}

The source is Paul D. Hanna's OEIS entry A393172 \cite{oeis172},
retrieved on 24 October 2026. Its initial sequence index is 1.
We also retain the formal zero constant coefficient at index zero.

All series are formal. The set $\N$ includes zero, and $[x^n]U$ denotes a
coefficient. A normalized series is $x+O(x^2)$; $O(x^r)$ means divisibility
by $x^r$, over any coefficient ring used below. We write $U\circ V=U(V(x))$
only when $V(0)=0$. Compositional inversion, denoted $\Rev(U)$ or $U\rev$,
is distinct from multiplicative inversion of a unit. Congruences are
coefficientwise; ordinary products and compositional iterates are distinct.

\section{Formal algebra used in the proof}

\begin{lemma}[Triangular construction]\label{lem:construction}
Let $K$ be a commutative ring. Suppose $\Phi$ preserves normalized series
and improves agreement by one degree:
\[
 G-H=O(x^r)\quad\Longrightarrow\quad
 \Phi(G)-\Phi(H)=O(x^{r+1})\qquad(r\ge2).
\]
Then $\Phi$ has a unique normalized fixed point. The same conclusion holds
for a family of series if the degree-$n$ equations, at every family index,
determine the degree-$n$ coefficients from strictly lower degrees.
\end{lemma}

\begin{proof}
The degree-$n$ coefficient of $\Phi(G)$ depends only on coefficients of
$G$ below degree $n$. Fix the constant and linear coefficients, and then
define every later coefficient successively. This gives a fixed point;
the same induction proves uniqueness. For a family, perform the induction
on degree simultaneously at every family index. The same induction applies with only the constant coefficient prescribed,
provided agreement is improved starting at degree one. It also applies to
an operator on all series when agreement is improved starting at degree zero.
For the operators used in this paper, the recurrence requires
only ring operations, so it also constructs integral solutions when
$K=\Z$. Equivalently, iterating $\Phi$ from any normalized initial series
fixes one further coefficient at each step.
\end{proof}

\begin{lemma}[Inversion and transport]\label{lem:inverse}
A normalized series over any commutative ring has a unique normalized
two-sided compositional inverse. A constant-one series has a unique
multiplicative inverse with constant coefficient one. These inverses and
valid composition commute with coefficient ring homomorphisms.
Normalized compositional inversion preserves agreement modulo $x^r$.
\end{lemma}

\begin{proof}
For normalized $G$, the degree-$n$ equation in $G(H)=x$ determines $h_n$
with multiplier one from earlier coefficients. Construct a right inverse
$J$ of $H$. Associativity gives
$G=G\circ(H\circ J)=(G\circ H)\circ J=J$, so $H\circ G=x$ as well.
The same associativity argument proves uniqueness. Multiplicative
inversion follows similarly from the coefficient recurrence for $UV=1$.

Each coefficient of a valid composition is a finite polynomial in the
coefficients of its arguments. Ring homomorphisms preserve these
polynomials and the defining equations, so uniqueness gives transport of
inversion. If $G-H=O(x^r)$, composition with any normalized series
preserves that agreement. Composing with the respective inverses then
shows $\Rev(G)-\Rev(H)=O(x^r)$.
\end{proof}

\begin{lemma}[Binary series and doubling]\label{lem:binary}
There is a unique normalized integral series $C=x+C^2$. Its reduction
modulo two is
\[
 P(x)=\sum_{k\ge0}x^{2^k}.
\]
For series $U,V$ over $\Z/4\Z$, the equality $2U=2V$ is equivalent to
equality of their reductions modulo two.
\end{lemma}

\begin{proof}
The operator $G\mapsto x+G^2$ improves agreement because
$G^2-H^2=(G-H)(G+H)$ and normalized $G+H$ has zero constant coefficient.
In characteristic two, $U(x)^2=U(x^2)$ for every formal series. Thus the
normalized equation $U=x+U^2$ forces the coefficient at one to be one,
every higher odd-index coefficient to be zero, and the coefficient at
$2m$ to equal the coefficient at $m$. This gives exactly $P$.

For scalar residues modulo four, $2u=2v$ holds precisely when $u$ and $v$
have the same parity. Apply this statement coefficientwise. In
particular, this argument does not cancel a factor of two in a ring in
which two is a zero divisor.
\end{proof}

\section{Source construction and the modulus-four theorem}

Let $C=x+C^2$ and put $N=x-x^2$. Then $N\circ C=x$, and
Lemma~\ref{lem:inverse} also gives $C\circ N=x$. We use $A$ for the source sequence A393172 \cite{oeis172}.

\begin{lemma}[A useful square equation]\label{lem:square}
For normalized $T,P\in K[[x]]$, there is a unique normalized solution of
\[
 G=T+(G\circ P)^2.
\]
The construction commutes with coefficient ring homomorphisms.
\end{lemma}

\begin{proof}
If $G-H=O(x^r)$, then $(G-H)\circ P=O(x^r)$. The sum
$G\circ P+H\circ P$ has zero constant coefficient, so the difference
of the two squares is $O(x^{r+1})$. Normalization is preserved.
Lemma~\ref{lem:construction} gives existence and uniqueness; mapping the
equation and using uniqueness gives transport.
\end{proof}

\begin{theorem}[A393172 modulo four]\label{thm:catalanfour}
The unique normalized integral solution of
\[
 A\circ N=x+A^2 \tag{1}
\]
has, for every $n>1$,
\[
 [x^n]A\equiv
 \begin{cases}
 2\pmod4,& n=2^k\text{ for some }k\ge0,\\
 0\pmod4,& \text{otherwise}.
 \end{cases}
\]
\end{theorem}

\begin{proof}
Composing (1) with $C$ gives
$A=C+(A\circ C)^2$. Conversely, composing this equation with $N$
gives (1). Lemma~\ref{lem:square} therefore proves integral existence
and uniqueness for the literal source equation.

Over $\Z/4\Z$, the normalized candidate $D=2C-x$ satisfies
\[
 D\circ N=2x-N=x+x^2,
 \qquad D^2=x^2+4(C^2-xC)=x^2.
\]
Thus $D$ solves (1), and uniqueness gives $A=2C-x$ modulo four. For
$n>1$, the coefficient of $x$ contributes nothing and Lemma~\ref{lem:binary}
shows that twice the coefficient of $C$ is two precisely at powers of
two, and zero elsewhere.
\end{proof}


\section{Separate Lean proof and checked scope}

The entry file is \nolinkurl{lean/Entries/A393172.lean}. Its explicit
theorem statements and proof terms use the namespace
\nolinkurl{FinalCollection.A393172}. The complete imported implementations
and their shared mathematical dependencies are included in the collection.
The following local theorem names give the proof map:

\begin{multicols}{2}\raggedright\begin{itemize}\setlength{\itemsep}{2pt}\setlength{\parsep}{0pt}
\item \nolinkurl{first_existence}.
\item \nolinkurl{first_definition}.
\item \nolinkurl{first_residue}.
\item \nolinkurl{first_arbitrary_source}.
\item \nolinkurl{complete_series_mod_four}.
\item \nolinkurl{divisibility_complement}.
\end{itemize}\end{multicols}

The main implementation modules are
\begin{quote}\raggedright
\noindent\nolinkurl{Round4.CatalanComposition}\par
\noindent\nolinkurl{Round4.SquareComposition}\par
\end{quote}
The project pins Lean 4.19.0 and Mathlib
\nolinkurl{c44e0c8ee63ca166450922a373c7409c5d26b00b}.
The supplied collection report records fresh source compilation and
transitive axiom audits. All coefficient conclusions are unbounded.

\paragraph{Limits and status.}
Normalized source equation $A(x-x^2)=x+A(x)^2$ is constructed uniquely. The standard Catalan parity mechanism is not claimed as a new general result.
No global novelty, exclusive priority, peer review, or OEIS acceptance is asserted.

\begin{thebibliography}{9}
\raggedright
\bibitem{oeis172}
Paul D. Hanna, \emph{OEIS A393172}, 4 February 2026.
\url{https://oeis.org/A393172}. Retrieved 24 October 2026.
\end{thebibliography}
\end{document}
